\chapter{绪论}

\label{cha:introduction}
\section{选题背景与意义}
\label{sec:background}

WebAssembly（Wasm） 是一种新兴的二进制指令集架构和代码分发格式，最初由 Mozilla、Google、Microsoft 等 Web 厂商联合提出，旨在解决 Web 平台中高性能计算能力不足、跨语言支持受限及代码可移植性差等问题\cite{haas2017bringing}。Wasm 被设计为一种架构无关、可验证且安全的低级中间表示，具备安全、高效、可移植等特点，能够作为多种高级语言的统一编译目标。

自提出以来，WebAssembly 便被广泛应用于 Web 场景，在浏览器中为 C/C++ 等高级语言提供接近原生性能的执行能力。与 JavaScript 相比，Wasm 采用静态类型和结构化控制流，能够更好地支持编译器优化和高效执行。目前，WebAssembly 现已成为继 HTML、CSS 和 JavaScript 后的第 4 个 Web 语言标准规范，并得到了所有主流浏览器的支持。C/C++、C\#、Go、Rust 等多种高级语言均已具备成熟的 Wasm 编译工具链，使复杂计算逻辑可以安全地运行在浏览器沙箱中\cite{mdn:wasm-c-guide,mdn_rust_wasm,gopherjs,zakai2011emscripten}。

随着 WebAssembly 技术的成熟，其应用范围逐渐从 Web 扩展到非 Web 场景。2019 年，Mozilla 推出了 WebAssembly System Interface（WASI），为 Wasm 程序提供了一套标准化的操作系统接口，如文件系统访问、时间管理和网络通信等\cite{wasi_standardizing}。WASI 的引入使 Wasm 不再依赖浏览器环境，而是能够作为一种轻量级、安全的可移植执行单元运行于主机、云平台和边缘设备之上。在此背景下，WebAssembly 被视为容器和虚拟机之外的一种新型运行时抽象，在微服务、边缘计算、物联网等领域得到广泛的探索和实践\cite{liu2021aerogel,sok2024webassembly}。

随着 Wasm 应用场景的不断拓展，其性能特性逐渐成为研究关注的重点之一。已有大量工作从不同角度对 WebAssembly 的性能进行了分析评估，例如比较不同语言编译为 Wasm 后的执行效率差异\cite{prabakar2024webassembly}，分析主流 Wasm 运行时（如 wasmtime、wasmer、wasm3、WAMR 等）的性能表现\cite{zhang2025research,dierickxcomparative}，以及对 Wasm 与 JavaScript 或原生代码之间的性能差距进行实验性研究等\cite{jangda2021not,sunarto2023systematic,yan2021understanding}。然而，现有研究大多基于运行时测量，侧重于性能结果的对比，而对性能差异产生原因的分析仍较为有限，尤其缺乏基于程序静态特征的系统性研究。目前尚未有工作在不依赖程序实际运行的前提下，从程序结构层面系统分析并初步判别其在 Wasm 与 native 执行模式下相对性能差异的成因与方向。

基于上述研究空白，本文提出了一种基于程序静态特征的方法，探索 WebAssembly 与原生代码之间性能差异的潜在原因。通过对程序中计算密集度、内存访问模式和控制流结构等静态特征进行建模，分析这些特征与性能表现之间的关联关系，为理解 WebAssembly 的性能边界提供新的视角。

本研究结果可为开发者进行技术选型、理解 Wasm 性能边界提供参考依据。具体而言，所构建的静态特征体系与分类模型可作为开发早期的粗粒度筛查工具，在无需实际运行程序的前提下，初步评估目标程序在 Wasm 环境下出现显著性能劣势的可能性；在云原生和边缘计算等场景中，开发者可结合静态分析结果，综合考虑 Wasm 在安全隔离、可移植性与执行效率方面的权衡，辅助制定部署策略。需要指出的是，静态分析结果应作为参考依据而非决策依据，最终选型仍需结合目标环境下的实测性能加以验证。


\section{本文的论文结构与章节安排}

\label{sec:arrangement}

本文共分为六章，各章节内容安排如下：

第一章为绪论。简单说明了本文章的选题背景与意义。

第二章综述了 WebAssembly 核心设计、相关性能研究以及静态特征分析等相关工作。

第三章介绍实验框架、基准程序设计、性能测量方法和数据集构建过程。

第四章给出实验结果与统计分析，重点讨论静态特征与性能差距之间的关系，并结合典型程序开展案例分析。此外还构建了一个二分类模型，对程序在 Wasm 环境下是否可能出现显著性能劣势进行粗粒度判别。

第五章对实验发现进行讨论，分析其与已有工作的联系、实践意义以及研究局限。

第六章总结全文工作，并展望未来可进行的进一步探索研究。